% SPDX-License-Identifier: CC-BY-4.0
% Copyright (c) 2026 Martin Schröder <info@swedishembedded.com>

\section{Discussion}\label{sec:discussion}

\paragraph{The exam, not the adapter, was the weakest link.}
The strongest conclusion from this work is about measurement. Every statistic the first
system produced was either uninformative (the anchor, whose bound sat inside its noise) or
about something other than the claim (the gate's improvement on paraphrases of trained
tasks, an effect of the prompt credited to the fine-tune). When those were corrected, the
remaining evidence on the adapter is a small positive difference over the prompted base with
an interval that, at the exam sizes available, includes zero. A power analysis would have
shown this before the first run: seven families, one vote each, cannot reject at
$\alpha=0.05$ unless every discordant family goes the same way and there are at least five of
them.

\paragraph{What the data support about the adapter.}
Three statements are supported. (i) A persona prompt alone raises the share of tasks judged
right and cuts invented specifics sharply (\cref{tab:exams}); a fine-tune must be compared
with it. (ii) Candidates trained on dialogues alone are indistinguishable from the prompted
base (v6: 14 against 14), which is what distilling a prompted teacher predicts
\citep{askell2021general}. (iii) The candidates that saw the author's own text word for word
are, in the pilot exams, never below the prompted base and above it by 1 to 8 tasks in all five
exams, a direction consistent with the closest comparable study
\citep{chakrabarty2025readers}; but the size is within what chance produces at
these sample sizes.

\paragraph{Why the voice view is the main lever, and why that is a hypothesis.}
Dialogues written by an 8B teacher under the persona prompt add little to what the prompt
already does. The author's own text is the only data that carries something the prompt cannot.
The first system gave this text little weight (210 fixed-prompt fragments, a loss that
favoured short teacher answers, a monitoring set that was dialogue only, a learning rate that
never decayed). Whether the corrected recipe with content-described chunks adds a measurable
amount is exactly what the pilots on the frozen exam are for; this paper does not claim
it.

\paragraph{Methodology that generalised.}
Four practices appeared repeatedly and are cheap: reserve the exam before generating
anything; fix the primary comparison before running it; store every per-task verdict so that
analyses can be redone offline and a power analysis can use the exam's own discordance; and
treat each instrument (judge, grader, split, bootstrap) as a thing to be measured, with
controls whose answer is known and a simulation of its false-positive rate.

\section{Threats to Validity}\label{sec:threats}

\paragraph{Construct validity.}
``Judged right'' means that a model judge says the answer gives what a source passage
says in reply to a generated task. That is a proxy for ``thinks like the author''; it rewards
knowledge of the passage and position, not prose style, and is silent on whether the reasoning is
the author's. A model that reads the exam's tasks and the passage would solve them, so the
tasks test the author's views as stated in the reserved letters. The invented-specifics
count is a code proxy for fabrication: it counts numbers and names absent from the source and
partly measures length.

\paragraph{Judge.}
Agreement with a person has not been measured; the judge was validated on controls only.
Judge and policy are the same model family, which invites self-preference that we cannot
exclude. The same family also wrote the tasks and the training data, so style and
topic distributions of exam and training tasks are related even though their source families
are disjoint.

\paragraph{Pretraining contamination.}
The letters are public domain and probably appear in the pretraining data of the base. The
exam is held out from fine-tuning, not from pretraining. Pretraining exposure would favour
all arms equally, including the prompted base, so it should not bias the primary comparison,
but it limits what ``held out'' means.

\paragraph{Adaptive reuse of the pilot exams.}
Recipes were developed over several runs examined on the same 31- and 35-task sets. Even if
an effect were observed there, selection on a small shared test set would overstate it. This is
the reason for the dev suite. The pilots on the frozen exam use its first 25 families by name, so
the final test is not unseen once they run; its analysis should report those 25 families and the
remaining ones separately.

\paragraph{Statistical validity.}
Seven-unit family-level tests have the power shown in \cref{fig:power}. A percentile
bootstrap over seven clusters is anti-conservative (false-positive rate 0.07 in simulation),
and we use it only on the larger exam. The power simulations assume discordance and ICC rather than
using values estimated from the data, and use 400 replicates. Pooled counts are descriptive only. There is one
training seed per recipe, so run-to-run variance of the adapter is unknown; one seed of the
v7 data has been trained (a second adapter) but not examined.

\paragraph{Internal validity.}
The recipe changed between candidates (data, learning rate, schedule, loss normalisation,
monitoring), so no pair of rows in \cref{tab:exams} isolates a single change. The earlier
candidates were trained at a learning rate (3e-4) different from the current default; v6,
ablation and adams did not record theirs.

\paragraph{External validity.}
Two authors, both eighteenth-century American, writing in English, whose corpora are large
(Jefferson) or modest (Adams), one base model, one adapter rank. Nothing here bears on other
languages, modern authors, or authors with little text.

\section{Limitations and Not Yet Established}\label{sec:limits}

\begin{itemize}
\item \textbf{No result on the frozen final test.} It has not been run in full. Pilots on its first 25
  families are listed in \cref{tab:pilot}; those not completed are not measured.
\item \textbf{No human labels.} Judge agreement, kappa, false-right and false-wrong rates and
  length bias on real answers are not measured. The tooling exists.
\item \textbf{Voice has no judge-free result.} The held-out text likelihood of each arm and
  the training-text overlap (8-, 13- and 20-word runs, longest shared run) are implemented but
  have no completed measurement. A fine-tuned model that quotes its training text is
  memorising, not thinking like the author; we cannot yet separate the two.
\item \textbf{No seeds.} A single seed per recipe; no estimate of run-to-run variance.
\item \textbf{The serve check was not measured} for any Qwen3 candidate: the serving binary
  was not available to the runs, and the check reports ``not measured'', never a pass.
\item \textbf{Retrieval is not evaluated.} The pipeline has retrieval-conditioned and
  abstention records and an exam arm with retrieved passages, but no measured run used them;
  all results are for a parametric model asked without retrieval. Retrieval itself was measured
  separately on 70 situations (semantic search found the source paragraph in the top ten for
  66\% of them, lexical for 30\%).
\item \textbf{Which correction mattered.} No training correction (token-normalised loss,
  annealing, weight decay, alpha, epochs) has been ablated. They are justified by audit and
  literature, not by an effect on the exam.
\item \textbf{Throughput.} Without packing, most of the compute on the P40s is padding, which
  bounds the data a run sees; a run over the whole corpus is expensive.
\item \textbf{Genuine ``thinking like'' cannot be proven.} The defensible claim is the
  observable one: on text excluded from training, positions, prose and abstentions are
  closer to the author's than the specified baselines'. We have not established that.
\end{itemize}

\section{Conclusion}

We built and examined a pipeline that learns from one sentence and an author's writings. Its
most useful products are methodological: a family-level split that survives overlapping
editions, an exam reserved before generation and sized by simulation, a judge with controls
and deterministic overrides, and a catalogue of defects that each changed a reported number.
Corrected, the evidence is that a persona prompt does much of the work, that no fine-tuned
candidate is statistically distinguishable from the prompted base at the sample sizes
measured, that the point estimates favour the adapters that saw the author's own text, and that
a properly powered test (39 families, 244 tasks) exists but has not yet been run in full. Until it is run, ``the adapter beats the prompt'' is a hypothesis.
